\scanpage{10}
\uline{Exercice} \handwritten{4}.\handwritten{10}.\handwritten{6} . Soient \(S\) un \(\underline{U}\)-site dont la topologie est moins fine que la canonique , \(\underline{\deleted{\unsure{S}}\handwritten{M}}\) une partie de \(\mathrm{Fl}(S)\) possédant les propriétés suivantes :

a) Les morphismes de \(\underline{M}\) sont quarrables (I \handwritten{10}.\handwritten{3}) et \(\underline{\unsure{\deleted{S}M|\deleted{M}S}}\) est stable par changement de base.

b) \(\underline{\deleted{S}M}\) contient les flèches identiques et est stable par composition .

c) Une flèche \(u:X\to Y\) telle qu'il existe une famille couvrante \(Y_i\to Y\) , avec les \(X\times_Y Y_i\to Y_i\) dans \(\underline{M}\) , est elle-même dans \(\underline{M}\) .

d) Pour tout objet \(X\) , toute famille couvrante de \(X\) est raffinée par une famille couvrante \((f_i:X_i\to X)\) où les \(f_i\) sont dans \(\underline{M}\).

Pour tout objet \(X\) , \deleted{considérons le site}\inserted{notons} \(S(X)\) (``petit site de \(X\)'')
% The typed insertion is in the LEFT margin; the two pen signs match.
\reversemarginpar
\margin{\figure{\handwritten{\begin{tikzpicture}[x=1ex,y=1ex,baseline=0ex]
\draw[line width=0.35pt] (0,-0.1)--(0.55,1.65);
\draw[line width=0.35pt] (0.65,2.2) ellipse (0.18 and 0.28);
\end{tikzpicture}}\inserted{le site}}}
\normalmarginpar
\handwritten{\begin{tikzpicture}[x=1ex,y=1ex,baseline=0ex]
\draw[line width=0.35pt] (0,0)--(0.65,1.3);
\draw[line width=0.35pt,rotate around={20:(0.79,1.43)}] (0.79,1.43) ellipse (0.23 and 0.18);
\end{tikzpicture}}
\deleted{dont la catégorie sous-jacente est}\inserted{constitué par} la sous-catégorie pleine de \(S/X\) formée des objets dont le morphisme structural est dans \(\underline{M}\) , munie de la topologie induite par celle de \(S\) (III \handwritten{3}.\handwritten{1}).

1) Montrer que les familles couvrantes \((f_i:X_i\to X)_{i\varepsilon I}\) où les \(f_i\) sont dans \(\underline{M}\) forment une prétopologie sur \(S\) (II \handwritten{1}.\handwritten{3}) qui engendre la topologie de \(S\) .

2) Pour toute flèche \(u:X\to Y\) de \(S\) , montrer que le foncteur changement de base par \(u\) est un foncteur continu (III \handwritten{1}.\handwritten{1}) de \(S(Y)\) dans \(S(X)\) . \unsure{\deleted{o}O|\deleted{O}o|O}n notera
\[
S(u)^{\handwritten{*}}:S(Y)^{\handwritten{\sim}}\to S(X)^{\handwritten{\sim}}
\]
le foncteur correspondant entre les catégorie\handwritten{s} de faisceaux (III \handwritten{1}.\handwritten{2},iv)).

3) Définir une équivalence entre le topos \(S^{\handwritten{\sim}}\) et la catégorie des systèmes
\[
(F_X)\ (X\varepsilon\operatorname{ob}S)\quad,\quad(\handwritten{\varphi_u})\ (u\varepsilon\mathrm{Fl}\ S)
\]
formés d'objets \(F_X\varepsilon\operatorname{ob}S(X)^{\handwritten{\sim}}\) , et pour toute flèche \(u:X\to Y\) dans \(S\) , d'un morphisme \(\handwritten{\varphi}_u:S(u)^{\handwritten{*}}(F_Y)\to F_X\) , ces systèmes
% archive sheet number: 9

